% ============================================================
% bidi-auto.tex  —  XeLaTeX template for plain-text mixed
% Latin / Hebrew input with ZERO markup in the source.
%
% The Unicode Bidirectional Algorithm (UBA) handles direction
% automatically. XeLaTeX + bidi respect it at the paragraph
% level. Font fallback via fontspec covers missing glyphs.
%
% Compile:
%   pandoc input.md -o output.pdf \
%     --pdf-engine=xelatex \
%     --template=bidi-auto.tex
%
% Or standalone:
%   xelatex bidi-auto.tex
% ============================================================

% \documentclass[12pt, a4paper]{article}

% \usepackage{iftex}
% \RequireXeTeX

% ------------------------------------------------------------
% SINGLE FONT covering Latin + Hebrew — the key to zero markup.
% Taamey Frank CLM or SBL BibLit both span the full range.
% fontspec will fall back to the FallbackFonts list for any
% glyph the main font lacks; order matters.
% ------------------------------------------------------------
% \usepackage{fontspec}

% \setmainfont{FreeSans}[
%   Ligatures    = TeX,
%   FallbackFonts = {SBL BibLit, SBL Hebrew, FreeSans, DejaVu Sans}
% ]

% That's it for fonts. No \newfontfamily, no \texthebrew{}.

% ------------------------------------------------------------
% BIDI — loads the Unicode Bidirectional Algorithm support.
% With bidi loaded, XeLaTeX reads the Unicode script properties
% of each character and sets run direction automatically.
% No \LR or \RL commands needed in the source.
% ------------------------------------------------------------
% \usepackage{fontspec}
% \defaultfontfeatures{Renderer=HarfBuzz}
% \setmainfont{Georgia}[
%   Ligatures    = TeX,
%   FallbackFonts = {Alef, Arial Hebrew, Source Sans 3,Lucida Grande, FreeSans}
% ]
% \usepackage{fontspec}
% \defaultfontfeatures{Renderer=HarfBuzz}

% \setmainfont{Avenir Next}

% \newfontfamily\hebrewfont{Arial Hebrew}[
%   Renderer = HarfBuzz,
%   Script   = Hebrew
% ]

% \usepackage{fontspec}
% \defaultfontfeatures{Renderer=HarfBuzz}

% \setmainfont{Avenir Next}

% \newfontfamily\hebrewfont{Alef}[
%   Renderer = HarfBuzz,
%   Script   = Hebrew
% ]

% % \usepackage[Hebrew]{ucharclasses}
% % \setTransitionsFor{Hebrew}{\hebrewfont}{\normalfont}
% \usepackage[Hebrew, IPAExtensions]{ucharclasses}

% \newfontfamily\ipafont{Source Sans 3}[Renderer=HarfBuzz]

% \setTransitionsFor{Hebrew}{\hebrewfont}{\normalfont}
% \setTransitionsFor{IPAExtensions}{\ipafont}{\normalfont}
% % \usepackage{ucharclasses}
% % \setTransitionsForHebrew{\hebrewfont}{\normalfont}
% \usepackage{bidi}
% ============================================================
% loani.tex  —  include-in-header snippet
% XeLaTeX: Georgia main, Alef Hebrew, Source Sans 3 IPA
% ucharclasses watches all three IPA Unicode blocks.
% bidi loaded BEFORE ucharclasses — required load order.
% ============================================================

\usepackage{iftex}
\RequireXeTeX

\usepackage{fontspec}
\defaultfontfeatures{Renderer=HarfBuzz}

\setmainfont{Georgia}[
  Ligatures     = TeX,
  FallbackFonts = {Alef, Arial Hebrew, Source Sans 3, Lucida Grande}
]

\newfontfamily\hebrewfont{Alef}[
  Renderer = HarfBuzz,
  Script   = Hebrew
]

\newfontfamily\ipafont{Source Sans 3}[
  Renderer = HarfBuzz
]

% ------------------------------------------------------------
% bidi FIRST — must be initialised before ucharclasses hooks
% ------------------------------------------------------------
\usepackage{bidi}

% ------------------------------------------------------------
% ucharclasses — all blocks that carry IPA characters:
%   IPAExtensions        U+0250–U+02AF
%   SpacingModifierLetters U+02B0–U+02FF  (ʰ ʷ ˈ ˌ ː etc.)
%   PhoneticExtensions   U+1D00–U+1D7F   (small caps, turned letters)
% ------------------------------------------------------------
\usepackage[
  Hebrew,
  IPAExtensions,
  SpacingModifierLetters,
  PhoneticExtensions
]{ucharclasses}

\setTransitionsFor{Hebrew}              {\hebrewfont}{\normalfont}
% \setTransitionsFor{IPAExtensions}        {\ipafont}{\normalfont\rmfamily}
\setTransitionsFor{SpacingModifierLetters}{\ipafont}{\normalfont\rmfamily}
\setTransitionsFor{PhoneticExtensions}   {\ipafont}{\normalfont\rmfamily}
\setTransitionsFor{IPAExtensions}{\color{red}\ipafont}{\normalfont\rmfamily}
% \setTransitionsFor{IPAExtensions}        {\bgroup\ipafont}{\egroup}
% \setTransitionsFor{SpacingModifierLetters}{\bgroup\ipafont}{\egroup}
% \setTransitionsFor{PhoneticExtensions}   {\bgroup\ipafont}{\egroup}
% \setTransitionsFor{IPAExtensions}       {\ipafont}   {\normalfont}
% \setTransitionsFor{SpacingModifierLetters}{\ipafont} {\normalfont}
% \setTransitionsFor{PhoneticExtensions}  {\ipafont}   {\normalfont}
% \setTransitionsFor{IPAExtensions}{%
%   \begingroup\ipafont
% }{%
%   \endgroup
% }
% \setTransitionsFor{SpacingModifierLetters}{%
%   \begingroup\ipafont
% }{%
%   \endgroup
% }
% \setTransitionsFor{PhoneticExtensions}{%
%   \begingroup\ipafont
% }{%
%   \endgroup
% }

\usepackage[hyperfootnotes=false]{hyperref}
% \renewcommand{\,}{\thinspace}
% \usepackage{etoolbox}
% \AtBeginEnvironment{footnote}{\hbadness=10000\hfuzz=10pt}
% \let\oldfootnote\footnote
% \renewcommand{\footnote}[1]{\oldfootnote{\LTR #1}}
% \usepackage{etoolbox}
% \patchcmd{\footnote}{\ignorespaces}{\leavevmode\LTR\ignorespaces}{}{}
% \let\bidi@saved@everypar\everypar
% \usepackage{endfloat}
% \AtEndDocument{\let\LTR\relax\let\RTL\relax}
% \AtBeginDocument{
%   \let\LTR\relax
%   \let\RTL\relax
% }
\usepackage[hyperfootnotes=false]{hyperref}
